\chapter{Perron--Frobenius Theory for Channels and Transfer Maps: Supporting Results}
\label{ch:qpf_supporting_results}

This appendix supplies the canonical-gauge calculations used in
Chapter~\ref{ch:qpf}.

\section{Canonical-gauge algebra}
\label{sec:app_qpf_gauge}

For the right-canonical gauge in
Theorem~\cite[``Right-canonical (unital) gauge'']{QICLean2026Blueprint}, each summand becomes
\begin{align}
    (S^{-1}A^iS)(S^{-1}A^iS)^\dagger
     & =S^{-1}A^iSS^\dagger(A^i)^\dagger(S^\dagger)^{-1}
    \notag                                               \\
     & =S^{-1}A^i\rho(A^i)^\dagger(S^\dagger)^{-1}.
    \notag
\end{align}
Hence
\begin{align}
    \sum_iA'^i(A'^i)^\dagger
     & =S^{-1}\E_A(\rho)(S^\dagger)^{-1}
    =S^{-1}\rho(S^\dagger)^{-1}
    =\Id.
    \notag
\end{align}
For the left-canonical gauge in Theorem~\ref{thm:left_canonical_gauge},
\begin{align}
    (SA^iS^{-1})^\dagger(SA^iS^{-1})
     & =(S^\dagger)^{-1}(A^i)^\dagger S^\dagger SA^iS^{-1}
    \notag                                                 \\
     & =(S^\dagger)^{-1}(A^i)^\dagger\sigma A^iS^{-1},
    \notag
\end{align}
and therefore
\begin{align}
    \sum_i(A'^i)^\dagger A'^i
     & =(S^\dagger)^{-1}
    \left(\sum_i(A^i)^\dagger\sigma A^i\right)S^{-1}
    =(S^\dagger)^{-1}\sigma S^{-1}
    =\Id.
    \notag
\end{align}
Finally, for the square-root gauge
(\ref{eq:qpf_tp_gauge}), self-adjointness of $\sigma^{1/2}$ gives
\begin{align}
    (B^i)^\dagger B^i
     & =\sigma^{-1/2}(A^i)^\dagger
    \sigma^{1/2}\sigma^{1/2}A^i\sigma^{-1/2}
    \notag                                                 \\
     & =\sigma^{-1/2}(A^i)^\dagger\sigma A^i\sigma^{-1/2}.
    \notag
\end{align}
Summing, substituting the adjoint fixed-point equation, and cancelling the
outer square-root factors yields
\begin{align}
    \sum_i(B^i)^\dagger B^i
     & =\sigma^{-1/2}
    \left(\sum_i(A^i)^\dagger\sigma A^i\right)\sigma^{-1/2}
    =\sigma^{-1/2}\sigma\sigma^{-1/2}
    =\Id.
    \notag
\end{align}


\begin{lemma}[Square-root trace-preserving gauge is a gauge equivalence]%
    \label{lem:tp_gauge_equiv}
    \lean{MPSTensor.gaugeEquiv_tpGauge}
    \leanok
    \uses{def:gauge_equiv}
    If $\sigma$ is positive definite and $B$ is defined by
    (\ref{eq:qpf_tp_gauge}), then $A$ and $B$ are gauge equivalent.
\end{lemma}

\begin{proof}\leanok
    The gauge matrix is $\sigma^{1/2}$, which is invertible because
    $\sigma$ is positive definite. Thus (\ref{eq:qpf_tp_gauge}) has the form
    of the gauge relation (\ref{eq:mps_gauge_equiv}).
\end{proof}
